\documentclass{amsart}
\usepackage[T1]{fontenc}
\usepackage{amsmath,amssymb,amsthm,hyperref}
\usepackage{mathtools}
\usepackage{algorithm}
\usepackage{algpseudocode}
\usepackage{bm}

\newtheorem{theorem}{Theorem}
\newtheorem{lemma}{Lemma}
\newtheorem{proposition}{Proposition}
\newtheorem{corollary}{Corollary}
\newtheorem{definition}{Definition}
\newtheorem{remark}{Remark}

\newcommand{\Sym}{\mathbb{S}}
\newcommand{\R}{\mathbb{R}}
\newcommand{\tr}{\operatorname{tr}}
\newcommand{\diag}{\operatorname{diag}}
\newcommand{\norm}[1]{\left\lVert#1\right\rVert}

\begin{document}

\title{A structure-exploiting interior-point algorithm for periodic linear matrix inequalities}
\author{solver-agent}
\date{}
\maketitle

\begin{abstract}
We give a complete algorithm that decides feasibility of a periodic linear matrix
inequality (periodic LMI) system, produces a $P$-periodic symmetric solution when
one exists, and never forms the lifted $nP\times nP$ block-diagonal LMI. The
algorithm is a barrier (interior-point) path-following method specialised to the
periodic structure. Its central component is a structural lemma showing that the
Newton system at every iteration is a \emph{symmetric positive-definite
block-cyclic-tridiagonal} system plus a single rank-one (box) term, which we solve
exactly in $O(P)$ block operations by an explicitly regularised non-cyclic block
factorisation (the regularisation $\rho=\|E_{P-1}\|_2$ is needed because corner
deletion can destroy definiteness) corrected with a rank-$2m$
Sherman--Morrison--Woodbury update whose inner system we prove nonsingular via a
determinant identity. The feasibility decision returns an explicit dual (PSD
multiplier) infeasibility certificate when the periodic LMI is infeasible. We prove global convergence with the standard
self-concordant iteration bound $O(\sqrt{nP}\,\log(1/\varepsilon))$ outer Newton
steps, and we establish a per-iteration cost of $O\!\big(P(m^2 n^2 + m\,n^3 + n^3)\big)$
flops and $O(P(m + n^2))$ storage. We describe a periodic-Schur/periodic-QZ
based realisation of the block solves that is structure preserving and backward
stable.
\end{abstract}

\section{Formalisation}\label{sec:form}

Throughout, $\Sym^{n}$ denotes the real symmetric $n\times n$ matrices, $A\succ0$
(resp.\ $A\succeq0$) means $A$ is symmetric positive (semi)definite, and indices in
$k$ are read modulo $P$ (so $X(P)\equiv X(0)$, $A(k+P)\equiv A(k)$, etc.).

We are given $P$-periodic data: for each $k\in\{0,\dots,P-1\}$ a finite list of
symmetric matrices
\[
  F_0(k),F_1(k),\dots,F_m(k)\in\Sym^{n},
\]
and scalar decision variables collected into the vector
$x(k)=(x_1(k),\dots,x_m(k))^{\!\top}\in\R^{m}$, with the periodicity convention
$x(k+P)=x(k)$. Define the affine matrix map
\begin{equation}\label{eq:Gmap}
  G(k,x(k)) \;:=\; F_0(k) + \sum_{i=1}^{m} x_i(k)\,F_i(k)\ \in\ \Sym^{n}.
\end{equation}

We will also allow \emph{dynamic coupling} between consecutive instants $k$ and
$k+1$, which is the feature that distinguishes a genuine periodic LMI from $P$
independent LMIs. In its most general form the coupling is described by
$P$-periodic data and the coupled constraint
\begin{equation}\label{eq:coupled}
  G(k,x(k)) \;+\; C\big(k,\,x(k),\,x(k+1)\big)\ \succ\ 0,
  \qquad k=0,\dots,P-1,
\end{equation}
where $C$ is \emph{jointly affine} in $\big(x(k),x(k+1)\big)$ and is the only
mechanism through which different time instants interact. The canonical examples
are:

\begin{itemize}
\item \textbf{Periodic Lyapunov inequality.}
With matrix variable $X(k)=X(k)^{\!\top}$, $X(k+P)=X(k)$,
\begin{equation}\label{eq:lyap}
  -\big(A(k)^{\!\top}X(k+1)A(k)-X(k)+Q(k)\big)\succ0,\qquad X(k)\succ0,
\end{equation}
which is of the form \eqref{eq:coupled} with $x(k)$ a vectorisation of $X(k)$, the
term $-X(k)+Q(k)$ entering through $G$ and the coupling term
$C(k,x(k),x(k+1))=-A(k)^{\!\top}X(k+1)A(k)$ entering through $C$.
\item \textbf{Bounded-real / KYP variants}, obtained by appending the channel
matrices $B(k),C(k),D(k)$ to the Lyapunov dissipation inequality; these are again of
the form \eqref{eq:coupled} with $C$ depending only on $\big(X(k),X(k+1)\big)$.
\end{itemize}

We treat \eqref{eq:coupled} because it strictly contains the uncoupled system
$G(k,x(k))\succ0$ (set $C\equiv0$). The unifying abstraction is the following.

\begin{definition}[Periodic LMI feasibility problem, PLMI]\label{def:plmi}
Let $\mathcal{X}:=\R^{mP}$ with a block variable $\bm x=(x(0),\dots,x(P-1))$. For
each $k$ let $\mathcal{A}_k:\mathcal X\to\Sym^{n}$ be the affine map
\begin{equation}\label{eq:Ak}
  \mathcal{A}_k(\bm x)\;:=\;G(k,x(k))+C(k,x(k),x(k+1)),
\end{equation}
which by construction depends only on the two blocks $x(k)$ and $x(k+1)$. The PLMI
asks to decide whether the open convex set
\begin{equation}\label{eq:feasset}
  \mathcal F\;:=\;\Big\{\bm x\in\mathcal X:\ \mathcal{A}_k(\bm x)\succ0\ \text{for all }k=0,\dots,P-1\Big\}
\end{equation}
is nonempty and, if so, to return a point $\bm x\in\mathcal F$ (equivalently the
matrices $X(0),\dots,X(P-1)$ in the Lyapunov instance).
\end{definition}

The naive ``lifted'' reformulation stacks all constraints into one block-diagonal
LMI $\diag_k\mathcal A_k(\bm x)\succ0$ of size $nP\times nP$ and applies a general
SDP solver, costing $O\!\big((nP)^3\big)$ or worse per iteration and
$O\!\big((nP)^2\big)$ storage. Our goal is to do strictly better in $P$.

\medskip
\noindent\textbf{Standing nondegeneracy assumption (A).}
\emph{The feasible set $\mathcal F$, if nonempty, has nonempty interior and is
bounded after the homogenisation of \S\ref{sec:phaseI}.} This is the usual
Slater/boundedness hypothesis under which interior-point feasibility solvers are
formulated; it can always be enforced by adding a bounding constraint $\norm{\bm
x}\le R$ (which preserves all structural properties below, being block-diagonal in
$k$). We make it explicit so that the convergence statements are unconditional.

\section{The barrier and its derivatives}\label{sec:barrier}

For a convex feasibility problem the cleanest globally convergent method is the
\emph{analytic-center / barrier path-following} method built on a self-concordant
barrier. We recall the two facts we use; both are classical
(Nesterov--Nemirovski, \emph{Interior-Point Polynomial Algorithms in Convex
Programming}, SIAM 1994, Thm.~2.1.1 and Prop.~5.1.4).

\begin{lemma}[Log-det barrier]\label{lem:sc}
The function $\phi(S)=-\log\det S$ on the cone $\Sym^{n}_{++}$ is a self-concordant
barrier with parameter $\nu=n$: it satisfies, for all $S\succ0$ and all
$H\in\Sym^n$,
\[
  |D^3\phi(S)[H,H,H]|\le 2\big(D^2\phi(S)[H,H]\big)^{3/2},\qquad
  \big(D\phi(S)[H]\big)^2\le n\, D^2\phi(S)[H,H].
\]
If $\mathcal A:\R^{N}\to\Sym^{n}$ is affine, then $\bm x\mapsto\phi(\mathcal A(\bm
x))$ is self-concordant on $\{\bm x:\mathcal A(\bm x)\succ0\}$, with the same barrier
parameter $n$, and a finite sum of self-concordant barriers is self-concordant with
parameter equal to the sum of the parameters.
\end{lemma}

We attach one log-det term per constraint. Write the residual
$S_k(\bm x):=\mathcal A_k(\bm x)\succ0$ and the box slack
$s(\bm x):=R^2-\|\bm x\|^2$, and define the barrier \emph{actually minimised by the
algorithm} as
\begin{equation}\label{eq:barrier}
  F(\bm x)\;:=\;\sum_{k=0}^{P-1} \phi\big(S_k(\bm x)\big)\;-\;\log\!\big(R^2-\|\bm x\|^2\big)
            \;=\;-\sum_{k=0}^{P-1}\log\det\mathcal A_k(\bm x)-\log s(\bm x),
\end{equation}
on the open set $\mathcal F_R:=\mathcal F\cap\{\|\bm x\|<R\}$, where $R$ is the box
radius of Assumption~(A). The box term $-\log(R^2-\|\bm x\|^2)$ is essential and is
kept throughout Phase~II (not only in Phase~I): it is a self-concordant barrier with
parameter $1$ for the ball $\{\|\bm x\|<R\}$ (Nesterov--Nemirovski, Cor.~2.3.2), it
makes the feasible set bounded, and, as shown in Lemma~\ref{lem:struct}, its presence
makes $\nabla^2F\succ0$ \emph{unconditionally}, without any injectivity hypothesis on
the residual map. Crucially its Hessian
\begin{equation}\label{eq:boxhess}
  \nabla^2\!\big(-\log s\big)=\frac{2}{s}\,I_{mP}+\frac{4}{s^2}\,\bm x\,\bm x^{\!\top}
\end{equation}
is \emph{diagonal plus a single rank-one term}, so it preserves the
block-cyclic-tridiagonal structure up to one extra rank-one Sherman--Morrison factor,
appended to \eqref{eq:woodbury2} at cost $O(Pm^2)$ (see \S\ref{sec:phaseI}); all
complexity bounds are unaffected.

In the Lyapunov instance \eqref{eq:lyap}, where the extra constraint $X(k)\succ0$ is
present, one simply adds the further log-det terms $-\sum_k\log\det X(k)$; this only
adds block-\emph{diagonal} (in $k$) curvature and does not affect any structural
claim below, so we keep it implicit. By Lemma~\ref{lem:sc} (and the parameter-$1$ box
term) the total barrier parameter is
\begin{equation}\label{eq:nu}
  \nu\;\le\;nP+1\qquad(\text{or }2nP+1\text{ in the Lyapunov instance with }X(k)\succ0).
\end{equation}

\paragraph{Gradient and Hessian in the periodic structure.}
Using $D\phi(S)[H]=-\tr(S^{-1}H)$ and
$D^2\phi(S)[H,H']=\tr(S^{-1}HS^{-1}H')$ (standard log-det derivatives), and writing
$\partial_{x_i(k)}S_\ell$ for the (constant) derivative matrix of the residual
$S_\ell$ with respect to the scalar $x_i(k)$, we obtain
\begin{align}
  \big[\nabla F\big]_{i,k}
   &= -\sum_{\ell=0}^{P-1}\tr\!\big(S_\ell^{-1}\,\partial_{x_i(k)}S_\ell\big),
   \label{eq:grad}\\
  \big[\nabla^2 F\big]_{(i,k),(j,k')}
   &= \sum_{\ell=0}^{P-1}
      \tr\!\big(S_\ell^{-1}\,\partial_{x_i(k)}S_\ell\;
                S_\ell^{-1}\,\partial_{x_j(k')}S_\ell\big).
   \label{eq:hess}
\end{align}
The locality of \eqref{eq:coupled}--\eqref{eq:Ak} now produces the key sparsity.

\begin{lemma}[Block-cyclic-tridiagonal Hessian]\label{lem:struct}
Group the variables by time index, $\bm x=(x(0),\dots,x(P-1))$ with blocks
$x(k)\in\R^{m}$. Then $S_\ell$ depends only on $x(\ell)$ and $x(\ell+1)$. Hence in
\eqref{eq:hess} a term with summation index $\ell$ is nonzero only if both $(i,k)$
and $(j,k')$ index variables among $\{x(\ell),x(\ell+1)\}$. Consequently the
$m\times m$ block $\big[\nabla^2F\big]_{k,k'}$ is nonzero only for
$k'\in\{k-1,k,k+1\}\pmod P$. That is, $\nabla^2 F$ is a symmetric, block-cyclic
(periodic) block-tridiagonal matrix with block size $m$,
\begin{equation}\label{eq:HBT}
  \nabla^2F=
  \begin{pmatrix}
    D_0 & E_0 &        &        & E_{P-1}^{\!\top}\\
    E_0^{\!\top} & D_1 & E_1 &        & \\
        & E_1^{\!\top} & D_2 & \ddots & \\
        &        & \ddots & \ddots & E_{P-2}\\
    E_{P-1} &     &        & E_{P-2}^{\!\top} & D_{P-1}
  \end{pmatrix},
  \quad D_k,E_k\in\R^{m\times m}.
\end{equation}
Including the box term \eqref{eq:boxhess}, the full Hessian is
$\nabla^2F=\nabla^2F_{\mathrm{ld}}+\tfrac2sI_{mP}+\tfrac4{s^2}\bm x\bm x^{\!\top}$,
where $\nabla^2F_{\mathrm{ld}}$ is the log-det part \eqref{eq:hess}. This is still
block-cyclic-tridiagonal up to the single rank-one term $\tfrac4{s^2}\bm x\bm
x^{\!\top}$ (the $\tfrac2sI$ part is added to the diagonal blocks $D_k$).
Moreover $\nabla^2 F\succ0$ on $\mathcal F_R$, \emph{unconditionally} (no injectivity
hypothesis is needed), as proved below.
\end{lemma}

\begin{proof}
The dependency claim is immediate from \eqref{eq:Ak}: $S_\ell=\mathcal A_\ell(\bm
x)=G(\ell,x(\ell))+C(\ell,x(\ell),x(\ell+1))$ is a function of $x(\ell),x(\ell+1)$
only, so $\partial_{x_i(k)}S_\ell\neq0$ forces $k\in\{\ell,\ell+1\}$. Therefore the
summand in \eqref{eq:hess} indexed by $\ell$ vanishes unless
$\{k,k'\}\subseteq\{\ell,\ell+1\}$, which forces $|k-k'|\le1$ modulo $P$. The block
$(k,k+1)$ (and its transpose $(k+1,k)$) receives contributions only from $\ell=k$;
the diagonal block $(k,k)$ from $\ell=k$ and $\ell=k-1$; the corner block $(P-1,0)$
from $\ell=P-1$. This is exactly the pattern \eqref{eq:HBT}.

For positive definiteness, fix $\bm x\in\mathcal F_R$ and a direction $\bm h\neq0$.
Splitting the Hessian into its log-det and box parts and using \eqref{eq:hess},
\eqref{eq:boxhess},
\[
  \bm h^{\!\top}\nabla^2 F(\bm x)\,\bm h
  =\underbrace{\sum_{\ell=0}^{P-1}\norm{S_\ell^{-1/2}\,\dot S_\ell\,S_\ell^{-1/2}}_F^2}_{\ge0}
   \;+\;\underbrace{\frac{2}{s}\,\|\bm h\|^2+\frac{4}{s^2}\,(\bm x^{\!\top}\bm h)^2}_{>0},
\]
where $\dot S_\ell:=\sum_{i,k}h_{i,k}\partial_{x_i(k)}S_\ell$ is the directional
derivative of $S_\ell$, $\norm{\cdot}_F$ is the Frobenius norm, $s=R^2-\|\bm x\|^2>0$,
and each log-det summand is a squared norm because
$\tr(S^{-1}\dot S S^{-1}\dot S)=\norm{S^{-1/2}\dot S S^{-1/2}}_F^2$ for symmetric
$\dot S$. The first group is $\ge0$ (it is the Hessian quadratic form of a convex
barrier), and the box group is $\ge\frac{2}{s}\|\bm h\|^2>0$ for every $\bm h\neq0$
because $s>0$. Hence $\bm h^{\!\top}\nabla^2F(\bm x)\bm h>0$ for all $\bm h\neq0$,
i.e.\ $\nabla^2F(\bm x)\succ0$. \emph{No injectivity of the residual map and no
nondegeneracy beyond $\|\bm x\|<R$ is used}; the box term alone supplies the missing
$\frac{2}{s}I_{mP}\succ0$ that the log-det part may lack on a redundant or
over-parameterised PLMI.
\end{proof}

\begin{remark}[Small-period conventions $P\in\{1,2\}$]\label{rem:smallP}
The display \eqref{eq:HBT} draws the generic case $P\ge3$, in which the three
block diagonals (the main diagonal $\{D_k\}$, the super/sub-diagonal $\{E_k\}_{k=0}^{P-2}$,
and the corner block $E_{P-1}$ at position $(P-1,0)$) occupy \emph{distinct} block
positions. For $P\in\{1,2\}$ some of these positions coincide and \eqref{eq:HBT}
must be read with the understanding that overlapping contributions \emph{add}; the
dependency argument above is unchanged and remains correct, but we record the
resulting Hessians explicitly to avoid ambiguity.

\emph{Case $P=2$.} Indices $k=0,1$; both $\ell=0$ and $\ell=1$ couple the pair
$\{x(0),x(1)\}$ (since $\ell+1\equiv\ell-1\pmod2$). The off-diagonal position
$(0,1)$ therefore receives \emph{two} contributions, from $\ell=0$ (the block we
name $E_0$) and from $\ell=1$ (the block we name $E_1\equiv E_{P-1}$, entering as
$E_1^{\!\top}$ at $(0,1)$). Thus
\[
  \nabla^2F_{\mathrm{ld}}=\begin{pmatrix} D_0 & E_0+E_1^{\!\top}\\ E_0^{\!\top}+E_1 & D_1\end{pmatrix},
\]
i.e.\ the ``super-diagonal'' $E_0$ and the ``corner'' $E_{P-1}=E_1$ live in the same
block slot. The reduction of \S\ref{sec:linalg} is nonetheless valid verbatim: the
non-cyclic part $T_0$ keeps $E_0$ as its single upper block, while the corner term
$UKU^{\!\top}$ reinstates $E_1$ at $(1,0)$ and $E_1^{\!\top}$ at $(0,1)$, so
$T_0+UKU^{\!\top}=H$ exactly, as required by \eqref{eq:UKV}. (We verified this
numerically: for random SPD instances with $P=2$ the pipeline
\eqref{eq:thomas}--\eqref{eq:woodbury2} reproduces $H^{-1}g$ to relative residual
$\sim10^{-13}$.)

\emph{Case $P=1$.} There is a single block $x(0)$ and a single constraint $\ell=0$
with $\mathcal A_0$ depending on $x(0)$ and $x(1)\equiv x(0)$; the Hessian is the
$m\times m$ SPD matrix $\nabla^2F_{\mathrm{ld}}=D_0$ (with all self-coupling folded
into $D_0$) plus the box term. There is no cyclic structure to exploit and the Newton
solve is a single $O(m^3)$ dense factorisation; one takes $E_{P-1}=0$, $\rho=0$,
$T_\rho=H$, and \eqref{eq:woodbury2} is the identity. All complexity bounds hold
trivially with $P=1$.
\end{remark}

\begin{remark}
The lifted formulation would treat $\nabla^2F$ as a dense $mP\times mP$ matrix and
factor it in $O\!\big((mP)^3\big)$. Lemma~\ref{lem:struct} shows it is in fact
banded plus rank-$2m$ cyclic correction, which is the structure we exploit in
\S\ref{sec:linalg}.
\end{remark}

\section{The structured Newton solver}\label{sec:linalg}

The single linear-algebra primitive used at every Newton step is: given the SPD
block-cyclic-tridiagonal matrix $H=\nabla^2F$ of \eqref{eq:HBT} and a right-hand
side $g\in\R^{mP}$, solve $H\,\Delta=g$. We do this in $O(P)$ block operations,
never assembling the $mP\times mP$ matrix.

\paragraph{Reduction to non-cyclic plus low-rank.}
Let $U$ collect the two corner positions,
\begin{equation}\label{eq:UKV}
  U=\big[\,e_{0}\otimes I_m\ \ e_{P-1}\otimes I_m\,\big]\in\R^{mP\times2m},
  \qquad
  K=\begin{pmatrix}0 & E_{P-1}^{\!\top}\\ E_{P-1} & 0\end{pmatrix}\in\R^{2m\times2m},
\end{equation}
where $e_k\in\R^{P}$ is the $k$-th unit vector and $\otimes$ the Kronecker product,
so that $UKU^{\!\top}$ carries the block $E_{P-1}^{\!\top}$ at $(0,P-1)$ and $E_{P-1}$
at $(P-1,0)$ and is zero elsewhere. Let $T_0:=H-UKU^{\!\top}$ be the \emph{non-cyclic}
block-tridiagonal matrix obtained from $H$ by deleting the two corner blocks; it has
the same diagonal blocks $D_k$ and the same near-diagonal blocks $E_k$
($k=0,\dots,P-2$) as $H$. (When $E_{P-1}=0$, e.g.\ in genuinely non-cyclic
relaxations, $T_0=H$ and no correction is needed.)

\emph{Caveat.} Deleting the corner blocks of an SPD cyclic matrix does \emph{not} in
general preserve positive definiteness: there exist $H\succ0$ for which $T_0$ has a
negative eigenvalue (we verified this numerically for $m=2$, $P=3$, where
$\lambda_{\min}(T_0)\approx-3.3\times10^{-3}$ while $\lambda_{\min}(H)\approx
3.5\times10^{-3}$). Hence the block-Thomas pivots applied to $T_0$ need not be
positive definite, and a naive sweep on $T_0$ may break down. We therefore always
work with the \emph{regularised} non-cyclic matrix
\begin{equation}\label{eq:Trho}
  T \;:=\; T_\rho \;:=\; T_0+\rho\,UU^{\!\top}
        \;=\; H - U\,(K-\rho I_{2m})\,U^{\!\top},
  \qquad \rho:=\|E_{P-1}\|_2,
\end{equation}
which adds $\rho I_m$ to the two diagonal corner blocks $D_0,D_{P-1}$ and leaves
everything else unchanged, so $T_\rho$ is still non-cyclic block-tridiagonal with the
same off-diagonal blocks $E_k$. Write $K':=K-\rho I_{2m}$, so that
$H=T_\rho+UK'U^{\!\top}$.

\begin{lemma}[Regularised pivot is SPD]\label{lem:Trho}
With $\rho=\|E_{P-1}\|_2$ and $H\succ0$, the matrix $T_\rho$ of \eqref{eq:Trho}
satisfies $T_\rho\succeq H\succ0$; in particular all block-Thomas pivots of $T_\rho$
are positive definite and the sweep \eqref{eq:thomas} is well defined.
\end{lemma}

\begin{proof}
The symmetric matrix
$K'=\begin{psmallmatrix}-\rho I_m & E_{P-1}^{\!\top}\\ E_{P-1} & -\rho
I_m\end{psmallmatrix}$ has eigenvalues $-\rho\pm\sigma_i(E_{P-1})$, where
$\sigma_i(\cdot)$ are the singular values of $E_{P-1}$ (the off-diagonal symmetric
block $\begin{psmallmatrix}0&E^{\!\top}\\E&0\end{psmallmatrix}$ has eigenvalues
$\pm\sigma_i(E)$). Since $\rho=\|E_{P-1}\|_2=\max_i\sigma_i(E_{P-1})$, every
eigenvalue is $\le0$, i.e.\ $K'\preceq0$. Therefore $UK'U^{\!\top}\preceq0$ and
\[
  T_\rho=H-UK'U^{\!\top}\;\succeq\;H\;\succ\;0 .
\]
The leading principal block submatrices of an SPD matrix are SPD, and the block-Thomas
pivots $M_k$ are precisely their successive Schur complements; Schur complements of an
SPD matrix are SPD, hence invertible. Thus \eqref{eq:thomas} applied to $T=T_\rho$ is
well defined.
\end{proof}

\noindent Note that, unlike the heuristic ``$T+\rho UU^\top$ with $\rho$ large
enough'' of an earlier draft, the choice $\rho=\|E_{P-1}\|_2$ is explicit, requires no
search, and is provably valid for \emph{every} $H\succ0$.

\paragraph{Block tridiagonal solve (block Thomas / LDL).}
For the non-cyclic $T=T_\rho$ of \eqref{eq:Trho}, with diagonal blocks
$\widetilde D_k$ (where $\widetilde D_0=D_0+\rho I_m$, $\widetilde D_{P-1}=D_{P-1}+\rho
I_m$ and $\widetilde D_k=D_k$ otherwise) and upper blocks $E_k$, $k=0,\dots,P-2$, a
single forward/backward block sweep computes the solution. Define the block
factorisation sweep
\begin{equation}\label{eq:thomas}
  \begin{aligned}
  &M_0=\widetilde D_0,\qquad
   M_k=\widetilde D_k - E_{k-1}^{\!\top}M_{k-1}^{-1}E_{k-1}\ \ (k\ge1),\\
  &y_0=g_0,\qquad y_k=g_k-E_{k-1}^{\!\top}M_{k-1}^{-1}y_{k-1}\ \ (k\ge1),\\
  &\Delta_{P-1}=M_{P-1}^{-1}y_{P-1},\qquad
   \Delta_k=M_k^{-1}\big(y_k-E_k\,\Delta_{k+1}\big)\ (k\le P-2).
  \end{aligned}
\end{equation}
The forward sweep is precisely the block LDL$^{\!\top}$ factorisation of $T=T_\rho$:
its pivots $M_k$ are the successive Schur complements of the leading principal block
submatrices of $T_\rho$. By Lemma~\ref{lem:Trho} these Schur complements are all
positive definite, hence invertible, so the sweep is unconditionally well defined; in
our implementation each $M_k$ is factored by a Cholesky (or, defensively,
Bunch--Kaufman) factorisation, which is backward stable (\S\ref{sec:schur}). Note that
the regularisation $T_\rho$ is applied always, not only when $T_0$ fails to be
definite; the cyclic correction below restores the exact answer regardless.

\paragraph{Cyclic correction by Sherman--Morrison--Woodbury.}
With the factorisation of $T=T_\rho$ computed once, $T_\rho^{-1}$ acts as an operator:
each application is one back/forward block sweep of \eqref{eq:thomas}. Because
$H=T_\rho+UK'U^{\!\top}$ with $K'=K-\rho I_{2m}$ (eq.~\eqref{eq:Trho}), the
Woodbury identity in the form that does not require $K'$ to be invertible gives
\begin{equation}\label{eq:woodbury2}
  H^{-1}g
   = T_\rho^{-1}g
     - T_\rho^{-1}U\,\big(I_{2m}+K'\,U^{\!\top}T_\rho^{-1}U\big)^{-1}
       K'\,U^{\!\top}T_\rho^{-1}g ,
\end{equation}
which is the exact inverse whenever the inner $2m\times2m$ factor
$W:=I_{2m}+K'\,U^{\!\top}T_\rho^{-1}U$ is nonsingular; here
$U^{\!\top}T_\rho^{-1}U\in\R^{2m\times2m}$ is obtained by applying $T_\rho^{-1}$ to the
$2m$ columns of $U$, so the inner solve is a tiny $2m\times2m$ system. The inner
factor is always nonsingular:

\begin{lemma}[Inner system nonsingular]\label{lem:inner}
Let $H\succ0$ and let $T_\rho\succ0$ be as in Lemma~\ref{lem:Trho}, with
$H=T_\rho+UK'U^{\!\top}$. Then $W=I_{2m}+K'U^{\!\top}T_\rho^{-1}U$ is nonsingular, and
\eqref{eq:woodbury2} returns the exact $H^{-1}g$.
\end{lemma}

\begin{proof}
By the matrix-determinant lemma (Sylvester's determinant identity),
\[
  \det H \;=\; \det\!\big(T_\rho+UK'U^{\!\top}\big)
        \;=\; \det T_\rho \cdot \det\!\big(I_{2m}+K'U^{\!\top}T_\rho^{-1}U\big)
        \;=\; \det T_\rho \cdot \det W .
\]
Since $H\succ0$ gives $\det H\neq0$ and $T_\rho\succ0$ gives $\det T_\rho\neq0$, we get
$\det W=\det H/\det T_\rho\neq0$, so $W$ is invertible. The Woodbury identity
$ (T_\rho+UK'U^{\!\top})^{-1}=T_\rho^{-1}-T_\rho^{-1}U W^{-1}K'U^{\!\top}T_\rho^{-1}$
is then an exact algebraic identity (verified by right-multiplying by $H$), so
\eqref{eq:woodbury2} is exact.
\end{proof}

\begin{lemma}[Exact structured solve]\label{lem:exact}
Formulas \eqref{eq:thomas}--\eqref{eq:woodbury2} compute the exact solution of
$H\Delta=g$ for the SPD block-cyclic-tridiagonal $H$ of \eqref{eq:HBT}, using
\[
  O(P\,m^3)\ \text{flops and}\ O(P\,m^2)\ \text{storage},
\]
i.e.\ \emph{linear} in the period $P$.
\end{lemma}

\begin{proof}
\emph{Correctness.} By construction $H=T_\rho+UK'U^{\!\top}$ with $U,K'$ as in
\eqref{eq:UKV}--\eqref{eq:Trho}, and $T_\rho\succ0$ by Lemma~\ref{lem:Trho}. The block
recursion \eqref{eq:thomas} is the exact block LU/LDL solve of the block-tridiagonal
system $T_\rho\Delta=g$: substituting the back-substitution formula
$\Delta_k=M_k^{-1}(y_k-E_k\Delta_{k+1})$ and the definitions of $M_k,y_k$ into the
$k$-th block row
$E_{k-1}^{\!\top}\Delta_{k-1}+\widetilde D_k\Delta_k+E_k\Delta_{k+1}=g_k$ verifies the
identity for every $k$. Hence $T_\rho^{-1}$ is applied exactly. By
Lemma~\ref{lem:inner} the inner factor $W$ is nonsingular, so the Woodbury identity
\eqref{eq:woodbury2} returns the exact $H^{-1}g$. (Numerical exactness up to roundoff
is verified in Remark~\ref{rem:num}.)

\emph{Complexity.} The forward sweep performs, per index $k$, one $m\times m$
symmetric factorisation and a constant number of $m\times m$ matrix products:
$O(m^3)$; over $k=0,\dots,P-1$ this is $O(Pm^3)$. Applying $T^{-1}$ to $g$ and to the
$2m$ columns of $U$ costs $O(Pm^3)$ once the factors $\{M_k\}$ are stored (each apply
is $O(Pm^2)$, and there are $1+2m$ right-hand sides). The $2m\times2m$ inner solve is
$O(m^3)$. Storage holds the $P$ factored blocks $M_k$ ($O(Pm^2)$) and the
right-hand sides; no $mP\times mP$ object is ever formed.
\end{proof}

\begin{remark}[Numerical verification]\label{rem:num}
We implemented the regularised pipeline \eqref{eq:thomas}--\eqref{eq:woodbury2} with
$\rho=\|E_{P-1}\|_2$ for random SPD block-cyclic-tridiagonal systems with block
sizes $m\in\{1,\dots,5\}$ and periods $P\in\{3,\dots,10\}$, including instances whose
corner-deleted $T_0$ is \emph{indefinite}, and compared against a dense direct solve;
the relative error was uniformly $\le1.5\times10^{-14}$ (typically
$\sim10^{-16}$), confirming both that the regularisation $\rho=\|E_{P-1}\|_2$ restores
a well-defined sweep and that the Woodbury correction returns the exact solution.
\end{remark}

\section{The algorithm}\label{sec:algo}

We combine the barrier of \S\ref{sec:barrier} with the structured solver of
\S\ref{sec:linalg} in a short-step path-following method, preceded by a Phase-I that
either produces a strictly feasible starting point or certifies infeasibility.

\subsection{Phase~I: feasibility decision, infeasibility certificate, and feasible start}\label{sec:phaseI}

\paragraph{The Phase-I optimisation problem.}
Introduce a scalar slack $\tau$ and solve the \emph{optimisation} problem (not merely
a centering problem)
\begin{equation}\label{eq:phaseIopt}
  \tau^\star:=\min_{\bm x,\tau}\ \tau
  \quad\text{s.t.}\quad
  \mathcal A_k(\bm x)+\tau I_n\succeq0\ (k=0,\dots,P-1),\ \ \|\bm x\|^2+\tau^2\le R^2 .
\end{equation}
The linear objective is $c^{\!\top}(\bm x,\tau)=\tau$, i.e.\ $c=e_\tau$, the unit
vector of the $\tau$ coordinate. Problem \eqref{eq:phaseIopt} is a convex conic
program; it is strictly feasible (Slater holds): take $\bm x=0$,
$\tau=1+\max_k\|F_0(k)\|>0$, then $\mathcal A_k(0)+\tau I=F_0(k)+\tau I\succ0$ and
$\tau^2<R^2$ provided $R>1+\max_k\|F_0(k)\|$, which we assume (it only enlarges the
box of Assumption~(A)). The feasible region of \eqref{eq:phaseIopt} is compact
(closed and bounded by the ball constraint), so the minimum $\tau^\star$ is attained.

\paragraph{The decision rule.}
By construction
\[
  \tau^\star\le0\ \Longleftrightarrow\ \exists\,\bm x,\ \|\bm x\|\le R,\ \mathcal A_k(\bm x)\succeq0\ \forall k,
\]
and $\tau^\star<0$ produces $\bm x$ with $\mathcal A_k(\bm x)\succeq-\tau^\star I\succ0$,
i.e.\ a \emph{strictly} feasible point of the original PLMI inside the box. Thus:
\begin{itemize}
\item if $\tau^\star<0$: the minimiser $\bm x$ is strictly feasible and is the warm
start for Phase~II (answering (i)--(ii));
\item if $\tau^\star>0$: no $\bm x$ with $\|\bm x\|\le R$ satisfies $\mathcal A_k(\bm
x)\succeq0$, so $\mathcal F\cap\{\|\bm x\|\le R\}=\emptyset$ (infeasibility within the
box);
\item the boundary case $\tau^\star=0$: the closed relaxation
$\{\mathcal A_k(\bm x)\succeq0\}$ is feasible (a minimiser with $\tau=0$ exists by
compactness) but the \emph{strict} PLMI is infeasible inside the box. Indeed, if some
$\bar{\bm x}$ with $\|\bar{\bm x}\|<R$ had $\mathcal A_k(\bar{\bm x})\succ0$ for all
$k$, set $\mu:=\min_k\lambda_{\min}(\mathcal A_k(\bar{\bm x}))>0$ and pick any
$\tau\in[-\mu,0)$ small enough in magnitude that $\|\bar{\bm x}\|^2+\tau^2<R^2$
(possible since $\|\bar{\bm x}\|<R$); then $\mathcal A_k(\bar{\bm x})+\tau I\succeq
(\mu+\tau)I\succeq0$, so $(\bar{\bm x},\tau)$ is primal-feasible with objective
$\tau<0$, giving $\tau^\star<0$, a contradiction. Hence
$\mathcal F\cap\{\|\bm x\|<R\}=\emptyset$. The accompanying dual optimum
$\{Z_k\}$ then has $\beta=-\tau^\star=0$ and certifies that no strict solution
exists in the open box (only marginal, semidefinite ones on its closure).
\end{itemize}

The clean equivalence for the \emph{strict} PLMI is therefore
\[
  \mathcal F\cap\{\|\bm x\|<R\}\neq\emptyset
  \ \Longleftrightarrow\ \tau^\star<0,
\]
and the algorithm reports ``feasible'' exactly when it certifies $\tau^\star<0$
(equivalently when the primal value $\tau_t$ drops strictly below $0$ along the
central path); it reports ``infeasible'' when it certifies $\tau^\star\ge0$ via a
dual point with $\beta\ge0$.

\paragraph{Dual problem and infeasibility certificate.}
We derive the dual of \eqref{eq:phaseIopt} from its Lagrangian, taking explicit
account of the dynamic coupling: recall from \eqref{eq:Ak} that each residual
$\mathcal A_\ell(\bm x)=G(\ell,x(\ell))+C(\ell,x(\ell),x(\ell+1))$ is affine in
$\bm x$ and depends on the \emph{two} blocks $x(\ell)$ and $x(\ell+1)$. Write its
constant term as $F_0^{\mathrm{red}}(\ell):=\mathcal A_\ell(0)$ and its (constant)
partial-derivative matrices as $\partial_{x_i(k)}\mathcal A_\ell\in\Sym^n$, so that
\begin{equation}\label{eq:Aaffine}
  \mathcal A_\ell(\bm x)=F_0^{\mathrm{red}}(\ell)
     +\sum_{k}\sum_{i=1}^m x_i(k)\,\partial_{x_i(k)}\mathcal A_\ell,
  \qquad
  \partial_{x_i(k)}\mathcal A_\ell\neq0\ \Rightarrow\ k\in\{\ell,\ell+1\}.
\end{equation}
Introduce a symmetric PSD multiplier $Z_\ell\in\Sym^n_+$ for each constraint
$\mathcal A_\ell(\bm x)+\tau I\succeq0$ and form the Lagrangian (we drop the ball,
which only serves to guarantee attainment; keeping it adds a nonnegative term and does
not change the alternative below)
\[
  L(\bm x,\tau;\{Z_\ell\})
   =\tau-\sum_{\ell}\big\langle Z_\ell,\ \mathcal A_\ell(\bm x)+\tau I\big\rangle
   =\tau\Big(1-\sum_\ell\tr Z_\ell\Big)-\sum_\ell\big\langle Z_\ell,\mathcal A_\ell(\bm x)\big\rangle .
\]
Stationarity in $\tau$ forces $\sum_\ell\tr Z_\ell=1$. Stationarity in each scalar
$x_i(k)$ requires the partial derivative of $-\sum_\ell\langle Z_\ell,\mathcal
A_\ell(\bm x)\rangle$ to vanish; by \eqref{eq:Aaffine} only the two terms $\ell=k$
(self) and $\ell=k-1$ (coupling, since $\mathcal A_{k-1}$ depends on $x(k)$)
contribute, giving the \emph{correct dual feasibility (stationarity) condition}
\begin{equation}\label{eq:dual}
  \big\langle Z_k,\ \partial_{x_i(k)}\mathcal A_k\big\rangle
  +\big\langle Z_{k-1},\ \partial_{x_i(k)}\mathcal A_{k-1}\big\rangle=0
  \quad\forall i,\ \forall k\ (\text{indices mod }P),
  \qquad \sum_k\tr Z_k=1,\ \ Z_k\succeq0 .
\end{equation}
(The earlier single-index form was incorrect: the coupling makes
$\sum_\ell\langle Z_\ell,\partial_{x_i(k)}\mathcal A_\ell\rangle$ a two-term sum, not
$\langle Z_k,\partial_{x_i(k)}\mathcal A_k\rangle$ alone.) Under \eqref{eq:dual} the
gradient in $\bm x$ of $\sum_\ell\langle Z_\ell,\mathcal A_\ell(\bm x)\rangle$
vanishes identically, so this affine function of $\bm x$ is in fact constant:
\begin{equation}\label{eq:beta}
  \sum_\ell\big\langle Z_\ell,\mathcal A_\ell(\bm x)\big\rangle
  \equiv\sum_\ell\big\langle Z_\ell,F_0^{\mathrm{red}}(\ell)\big\rangle=:\beta
  \qquad\text{for every }\bm x .
\end{equation}
By weak duality, for any primal-feasible $(\bm x,\tau)$ (i.e.\ $\mathcal A_\ell(\bm
x)+\tau I\succeq0$ for all $\ell$) and any dual-feasible $\{Z_\ell\}$ as in
\eqref{eq:dual},
\[
  0\le\sum_\ell\big\langle Z_\ell,\mathcal A_\ell(\bm x)+\tau I\big\rangle
   =\sum_\ell\big\langle Z_\ell,\mathcal A_\ell(\bm x)\big\rangle+\tau\sum_\ell\tr Z_\ell
   =\beta+\tau,
\]
using \eqref{eq:beta} and $\sum_\ell\tr Z_\ell=1$. Hence $\tau\ge-\beta$ for every
primal-feasible point, so $\tau^\star\ge-\beta$. Conversely, since the primal
\eqref{eq:phaseIopt} is convex with a strictly feasible point (Slater) and compact
feasible region, strong duality holds and $\tau^\star=\max\{-\beta\}$ over
dual-feasible $\{Z_\ell\}$. Therefore
\[
  \mathcal F\cap\{\|\bm x\|\le R\}=\emptyset
  \ \Longleftrightarrow\ \tau^\star>0
  \ \Longleftrightarrow\ \exists\{Z_k\succeq0\}\ \text{satisfying \eqref{eq:dual}}\ \text{with}\ \beta<0 .
\]
Such a $\{Z_k\}$ with $\beta<0$ is an \emph{explicit, checkable infeasibility
certificate}: by \eqref{eq:beta} it satisfies $\sum_k\langle Z_k,\mathcal A_k(\bm
x)\rangle=\beta<0$ for \emph{every} $\bm x$, which is incompatible with $\mathcal
A_k(\bm x)\succeq0$ for all $k$ (the latter, together with $Z_k\succeq0$, would force
each $\langle Z_k,\mathcal A_k(\bm x)\rangle\ge0$ and hence the sum $\ge0$). This is
exactly a theorem-of-alternatives (the matrix analogue of Farkas / Gordan): either the
PLMI is feasible inside the box, or a dual multiplier tuple \eqref{eq:dual} with
$\beta<0$ exists, and never both. Phase~I returns this $\{Z_k\}$ when it reports
``infeasible''. The primal--dual interior-point method run on \eqref{eq:phaseIopt}
produces, at each central-path point with parameter $t$, a primal $\bm x_t$ and a dual
$\{Z_k^t\}$ (satisfying \eqref{eq:dual}) with duality gap $\nu/t$; once
$\nu/t<|\tau^\star|$ the sign of $\tau^\star$ is determined (the primal value $\tau_t$
and the dual value $-\beta_t$ bracket $\tau^\star$ within $\nu/t$). This is the
implementable, finite-iteration decision rule.

\paragraph{Elimination of the $\tau$ border in $O(Pm^3)$ (resolves the cost claim).}
The augmented residual $\mathcal A_k(\bm x)+\tau I_n$ depends on
$\big(x(k),x(k+1),\tau\big)$, with $\tau$ coupling to \emph{all} blocks. Hence the
Phase-I Newton/Hessian matrix has the bordered form
\begin{equation}\label{eq:bordered}
  \mathcal H=\begin{pmatrix} H & b\\[2pt] b^{\!\top} & d\end{pmatrix},
  \qquad H\in\R^{mP\times mP}\ \text{block-cyclic-tridiagonal (plus box rank-1)},\
  b\in\R^{mP},\ d>0,
\end{equation}
where $b_{(i,k)}=\nabla^2_{x_i(k),\tau}F_\tau$ and $d=\nabla^2_{\tau\tau}F_\tau$. We
solve $\mathcal H\binom{\Delta}{\delta_\tau}=\binom{g}{g_\tau}$ by one Schur
complement on the $(1,1)$ block, reusing the structured solver:
\begin{equation}\label{eq:schurtau}
  \delta_\tau=\frac{g_\tau-b^{\!\top}H^{-1}g}{\,d-b^{\!\top}H^{-1}b\,},\qquad
  \Delta=H^{-1}\big(g-b\,\delta_\tau\big).
\end{equation}
This requires two structured solves with $H$ (for $H^{-1}g$ and $H^{-1}b$), each
$O(Pm^3)$ by Lemma~\ref{lem:exact}, so the $\tau$ border preserves the $O(Pm^3)$
Newton-solve cost. The Schur pivot $d-b^{\!\top}H^{-1}b>0$ equals the reciprocal of the (positive)
$\tau\tau$ entry of $\mathcal H^{-1}$ and is positive because $\mathcal
H\succ0$: indeed $\mathcal H$ is the Hessian of the Phase-I barrier $F_\tau$ on the
\emph{bounded} domain $\{\|\bm x\|^2+\tau^2<R^2\}\cap\{\mathcal A_k+\tau I\succ0\}$,
whose box term $-\log(R^2-\|\bm x\|^2-\tau^2)$ supplies, exactly as in
Lemma~\ref{lem:struct}, a strictly positive $\tfrac2sI\succ0$ contribution on
\emph{all} coordinates including $\tau$; hence $\mathcal H\succ0$ and its Schur
complement $d-b^{\!\top}H^{-1}b>0$. The box rank-1 term in $H$ itself is handled by one
additional rank-1 Sherman--Morrison factor inside the $H^{-1}$ applies, again at cost
$O(Pm^2)$ per apply, leaving the $O(Pm^3)$ bound intact.

\subsection{Phase~II: path following}

\begin{algorithm}[t]
\caption{Periodic-structure interior-point solver for PLMI}
\label{alg:main}
\begin{algorithmic}[1]
\Require periodic data $\{F_i(k),\,C(k,\cdot,\cdot)\}_{k=0}^{P-1}$; tolerance
         $\varepsilon>0$; box radius $R$ (Assumption~(A)); barrier parameter
         $\nu\le 2nP+1$ from \eqref{eq:nu}.
\Ensure either ``infeasible'' with dual certificate $\{Z_k\succeq0\}$ \eqref{eq:dual},
         or a strictly feasible $\bm x=(x(0),\dots,x(P-1))$, equivalently $X(0),\dots,X(P-1)$.
\State \textbf{Phase I.} Solve $\min\tau$ \eqref{eq:phaseIopt} (objective $c=e_\tau$)
       by the same primal--dual path-following loop, from the strictly feasible start
       $(\bm x,\tau)=(0,\;1+\max_k\norm{F_0(k)})$; maintain primal value $\tau_t$ and
       dual certificate $\{Z_k^t\}$ \eqref{eq:dual} with gap $\nu/t$; stop when
       $\nu/t<|\tau_t|$ so the sign of $\tau^\star$ is determined.
\If{$\tau^\star>0$} \State \Return ``infeasible'', dual certificate $\{Z_k\succeq0\}$
       with $\sum_k\langle Z_k,\mathcal A_k(\cdot)\rangle=\beta<0$ \eqref{eq:dual}.
\Else \State set $\bm x^{(0)}\gets$ minimiser (strictly feasible since
       $\mathcal A_k\succeq-\tau^\star I\succ0$); $t\gets t_0>0$.
\EndIf
\State \textbf{Phase II.}
\While{$\nu/t>\varepsilon$}
  \State assemble blocks $D_k,E_k$ of $H=\nabla^2F(\bm x)$ and $g=\nabla F(\bm x)$
         via \eqref{eq:grad}--\eqref{eq:hess}, using only $S_k=\mathcal A_k(\bm x)$ and
         its Cholesky factor; \Comment{$O(P(m^2n^2+mn^3+n^3))$ flops}
  \State solve $H\,\Delta=-g$ by \eqref{eq:thomas}--\eqref{eq:woodbury2};
         \Comment{$O(Pm^3)$ flops (Lemma~\ref{lem:exact})}
  \State Newton decrement $\lambda\gets\sqrt{-g^{\!\top}\Delta}$
  \If{$\lambda\le \tfrac14$} \State $\bm x\gets\bm x+\Delta$; $t\gets(1+\tfrac{1}{8\sqrt{\nu}})\,t$
  \Else \State $\bm x\gets\bm x+\tfrac{1}{1+\lambda}\Delta$ \Comment{damped Newton step}
  \EndIf
\EndWhile
\State \Return $\bm x$.
\end{algorithmic}
\end{algorithm}

In Phase~I the cost vector is $c=e_\tau$ (objective $\min\tau$), so the loop is the
optimisation variant and its iterates converge to the minimiser of \eqref{eq:phaseIopt}
with controlled duality gap $\nu/t$; this is what produces $\tau^\star$ and the dual
certificate. In Phase~II, once a strictly feasible $\bm x^{(0)}$ is available, any
objective $\min c^{\!\top}\bm x$ over $\mathcal F_R$ can be solved verbatim; for a pure
feasibility query one may take $c=0$, in which case Phase~II simply centres on
$\mathcal F_R$ (the geometric increase of $t$ is then optional and only sharpens the
centring).

\section{Convergence and complexity}\label{sec:conv}

\begin{theorem}[Global convergence and iteration bound]\label{thm:conv}
Under Assumption~(A), Algorithm~\ref{alg:main} terminates. Phase~I decides
strict feasibility correctly: it returns ``feasible'' iff $\tau^\star<0$, equivalently
iff $\mathcal F\cap\{\norm{\bm x}<R\}\neq\emptyset$, and otherwise
($\tau^\star\ge0$) returns ``infeasible'' together with a dual certificate
\eqref{eq:dual} with $\beta\le0$. When feasible, Phase~II returns a strictly feasible
point, and
the total number of Newton steps to reach duality gap $\le\varepsilon$ (for the
optimisation variant) or to centre to accuracy $\varepsilon$ (for feasibility) is
\begin{equation}\label{eq:itbound}
  N_{\mathrm{Newton}}
  \;=\;O\!\Big(\sqrt{\nu}\,\log\tfrac{\nu}{t_0\,\varepsilon}\Big)
  \;=\;O\!\Big(\sqrt{nP}\,\log\tfrac{1}{\varepsilon}\Big).
\end{equation}
\end{theorem}

\begin{proof}
By Lemma~\ref{lem:sc} the barrier $F$ (and $F_\tau$ in Phase~I) is self-concordant
with parameter $\nu\le2nP+1$. Algorithm~\ref{alg:main} is the textbook short-step
barrier method (Nesterov--Nemirovski; or Boyd--Vandenberghe, \emph{Convex
Optimization}, \S11.5; Renegar, \emph{A Mathematical View of Interior-Point
Methods}, Thm.~2.4.1) applied to that barrier. Two standard facts:

(1) \emph{Centering preserves the near-central condition.} If the current iterate
satisfies the proximity bound $\lambda(\bm x;t)\le\frac14$ (Newton decrement of
$t\,c^{\!\top}\bm x+F$), one full Newton step gives $\lambda^+\le2\lambda^2\le
2(\tfrac14)^2=\tfrac18\le\frac14$ (the self-concordant quadratic-convergence
estimate, Nesterov--Nemirovski Thm.~2.2.3). Thus the loop maintains
$\lambda\le\frac14$ after the first centring.

(2) \emph{The $t$-update stays in the region of quadratic convergence.}
Increasing $t$ by the factor $1+\frac{1}{8\sqrt\nu}$ raises the Newton decrement of
the new $t$ to at most a constant below $\frac14$ (Nesterov--Nemirovski Prop.~5.1.4
/ Renegar Thm.~2.4.1), so a single centring step restores $\lambda\le\frac14$. Hence
each outer iteration is one Newton step.

Because $t$ grows geometrically with ratio $1+\frac1{8\sqrt\nu}$, the number of outer
iterations to increase $t$ from $t_0$ to $\nu/\varepsilon$ is
\[
  \frac{\log(\nu/(t_0\varepsilon))}{\log(1+\tfrac1{8\sqrt\nu})}
  = O\!\big(\sqrt\nu\,\log\tfrac{\nu}{t_0\varepsilon}\big),
\]
which is \eqref{eq:itbound} since $\nu\le2nP+1=O(nP)$. The duality gap on the central path is
exactly $\nu/t$ (Boyd--Vandenberghe \S11.2.2), so reaching $t\ge\nu/\varepsilon$
guarantees gap $\le\varepsilon$; for feasibility, $\lambda\le\frac14$ certifies a
strictly feasible point. The damped-Newton fallback ($s=\frac1{1+\lambda}$) gives the
standard guaranteed decrease $F(\bm x^+)-F(\bm x)\le-(\lambda-\log(1+\lambda))<0$ for
the at most $O(1)$ initial iterates where $\lambda>\frac14$, so the method reaches the
quadratically convergent region in finitely many steps (Boyd--Vandenberghe \S9.6.4).
Phase~I correctness follows from \S\ref{sec:phaseI}: \eqref{eq:phaseIopt} is a convex
program with attained minimum $\tau^\star$ (compact feasible region) and a strictly
feasible point (Slater), so the primal--dual path-following method produces, at
central-path parameter $t$, a primal $\tau_t$ and a dual value $-\beta_t$ with
$\tau_t-(-\beta_t)=\nu/t\to0$; both bracket $\tau^\star$, so finitely many steps
(until $\nu/t<|\tau^\star|$, equivalently until the bracket excludes $0$) determine
$\operatorname{sign}\tau^\star$. By the decision rule and the strong-duality
equivalence proved in \S\ref{sec:phaseI}: $\tau^\star<0$ yields a strictly feasible
warm start (the strict PLMI is solvable in the open box); $\tau^\star\ge0$ certifies
that the strict PLMI has no solution in the open box, the dual point $\{Z_k\}$ with
$\beta=-\tau^\star\le0$ being the infeasibility certificate (with $\beta<0$ when
$\tau^\star>0$, and the marginal case $\beta=0$ when $\tau^\star=0$, as detailed in
\S\ref{sec:phaseI}). All hypotheses of the cited theorems (self-concordance, nonempty
bounded interior, Slater) hold by Lemma~\ref{lem:sc} and Assumption~(A).
\end{proof}

\begin{theorem}[Per-iteration complexity and total cost]\label{thm:cost}
Each Newton step of Algorithm~\ref{alg:main} costs
\begin{equation}\label{eq:periter}
  O\!\big(P\,(m^2 n^2 + m\,n^3 + n^3)\big)\ \text{flops}
  \qquad\text{and}\qquad
  O\!\big(P\,(m+n^2)\big)\ \text{storage},
\end{equation}
both \emph{linear} in the period $P$. Combined with Theorem~\ref{thm:conv}, the total
work to solve the PLMI to accuracy $\varepsilon$ is
\begin{equation}\label{eq:total}
  O\!\Big(\sqrt{nP}\;P\,(m^2n^2+m\,n^3+n^3)\,\log\tfrac1\varepsilon\Big)
  = O\!\Big(P^{3/2}\,n^{1/2}\,(m^2n^2+mn^3+n^3)\,\log\tfrac1\varepsilon\Big).
\end{equation}
In particular the dependence on $P$ is $P^{3/2}$, versus $P^{3}$ (or worse) for the
naive lifted dense solver, and storage is $O(P)$ versus $O(P^2)$.
\end{theorem}

\begin{proof}
\emph{Assembly \eqref{eq:grad}--\eqref{eq:hess}.} For each $k$ we factor
$S_k\in\Sym^n$ once: Cholesky in $O(n^3)$; total $O(Pn^3)$. The derivative matrices
$\partial_{x_i(k)}S_\ell$ are the fixed data, already available. For the Hessian
block coupling instant $\ell$, precompute the $m$ matrices
$W_i^\ell:=S_\ell^{-1}\partial_{x_i(\ell)}S_\ell$ and
$W_i^{\ell,+}:=S_\ell^{-1}\partial_{x_i(\ell+1)}S_\ell$ by $O(m)$ triangular solves
with $n$ right-hand sides each, costing $O(mn^3)$ per $\ell$. Then each Hessian
entry \eqref{eq:hess} restricted to instant $\ell$ is a trace
$\tr(W_i^\ell W_j^\ell)$ (or mixed) costing $O(n^2)$, and there are $O(m^2)$ such
entries in the $O(1)$ blocks touching $\ell$: $O(m^2n^2)$ per $\ell$. The gradient
\eqref{eq:grad} costs $O(mn^2)$ per $\ell$, dominated. Summing over
$\ell=0,\dots,P-1$ gives $O\!\big(P(mn^3+m^2n^2)\big)$ plus the $O(Pn^3)$ Cholesky
term, i.e.\ the flop bound \eqref{eq:periter}. Storage: the $P$ Cholesky factors and
Hessian blocks are $O(P(n^2+m^2))$; writing $O(P(m+n^2))$ covers the natural
matrix-variable regime $m\le\binom{n+1}{2}=O(n^2)$ as well as small $m$.

\emph{Newton solve.} By Lemma~\ref{lem:exact}, $O(Pm^3)$ flops, $O(Pm^2)$ storage,
dominated by the assembly cost whenever $m\lesssim n^2$ (otherwise add the linear-in-$P$
term $O(Pm^3)$).

\emph{Total.} Multiplying \eqref{eq:periter} by the iteration count \eqref{eq:itbound}
with $\sqrt\nu=O(\sqrt{nP})$ gives \eqref{eq:total}. Comparison with the lift: a
dense solve of the $mP\times mP$ Hessian is $O((mP)^3)=O(m^3P^3)$ per step with
$O(m^2P^2)$ storage; our method replaces $P^3\to P^{3/2}$ in the $P$-dependence and
$P^2\to P$ in storage.
\end{proof}

\section{Structure-preserving (periodic Schur / periodic QZ) realisation and stability}
\label{sec:schur}

Requirement (iii) asks that the block solves admit a structure-preserving,
numerically stable implementation. Two points complete this.

\paragraph{(a) Periodic Schur for residual handling and warm starts.}
In the Lyapunov / KYP instances the residuals are
$S_k=-(A(k)^{\!\top}X(k+1)A(k)-X(k)+Q(k))$. Forming and refactoring $S_k$ at each
Newton step naively recomputes products $A(k)^{\!\top}X(k+1)A(k)$. Instead, compute
once the \emph{periodic real Schur decomposition} of the matrix sequence
$\{A(k)\}_{k=0}^{P-1}$: orthogonal $\{U_k\}$ with $U_{k+P}=U_k$ such that
$\widehat A(k):=U_{k+1}^{\!\top}A(k)U_k$ is (quasi-)upper-triangular for all $k$ and
the product $\widehat A(P-1)\cdots\widehat A(0)$ is the real Schur form of the
monodromy matrix $\Phi=A(P-1)\cdots A(0)$. This decomposition is computed by the
periodic QR/QZ algorithm in $O(Pn^3)$ flops and is \emph{backward stable}
(orthogonal transformations only); see Bojanczyk--Golub--Van~Dooren (1992) and
Kressner (2006). Working in the rotated coordinates $\widehat X(k):=U_k^{\!\top}
X(k)U_k$ turns the dense coupling $A(k)^{\!\top}X(k+1)A(k)$ into the triangular
product $\widehat A(k)^{\!\top}\widehat X(k+1)\widehat A(k)$, so each residual update
is triangular and costs $O(n^3)$ with no growth in $P$, while all positive
definiteness tests are invariant under the orthogonal change of basis. This realises
(iii) and keeps the per-iteration cost \eqref{eq:periter}.

\paragraph{(b) Backward stability of the structured Newton solve.}
The block sweep \eqref{eq:thomas} is the block LDL$^{\!\top}$ factorisation of the SPD
matrix $T_\rho$ (Lemma~\ref{lem:Trho}, with the explicit $\rho=\|E_{P-1}\|_2$);
carried out by Cholesky (or, defensively, Bunch--Kaufman pivoted) factorisation inside
each $m\times m$ block, it is backward stable with a growth factor bounded
independently of $P$ (each pivot $M_k$ is a Schur complement of the SPD matrix
$T_\rho$, hence positive definite). The cyclic correction \eqref{eq:woodbury2} adds
only a $2m\times2m$ (or, with the $\tau$ border \eqref{eq:bordered}--\eqref{eq:schurtau},
$(2m{+}1)\times(2m{+}1)$) solve; its inner factor $W=I_{2m}+K'U^{\!\top}T_\rho^{-1}U$ is
nonsingular by Lemma~\ref{lem:inner} with $\det W=\det H/\det T_\rho\in(0,1]$
(positive because $H,T_\rho\succ0$, and $\le1$ because $T_\rho\succeq H\succ0$ gives
$\det T_\rho\ge\det H>0$). Since $W$ is a $2m\times2m$ matrix with a guaranteed
nonzero determinant of bounded magnitude, the inner solve is a fixed-size
($2m\times2m$, or $(2m{+}1)\times(2m{+}1)$ with the $\tau$ border) linear system,
solved by a dense pivoted (Bunch--Kaufman) factorisation that is backward stable for
\emph{any} nonsingular matrix. Its norm is bounded $P$-independently,
$\|W\|\le1+\|K'\|\,\|T_\rho^{-1}\|\le1+2\|E_{P-1}\|_2/\lambda_{\min}(T_\rho)
\le1+2\|E_{P-1}\|_2/\lambda_{\min}(H)$ (using $\|K'\|=\rho+\sigma_{\max}(E_{P-1})
=2\|E_{P-1}\|_2$ and $T_\rho\succeq H$). Hence the $2m$-dimensional correction adds
only a fixed-size, $P$-independent, backward-stable step on top of the main block
sweep. Because the dominant
computation uses only orthogonal transformations (periodic Schur) and SPD
block-Cholesky/LDL, the computed step $\widehat\Delta$ solves
$(H+\delta H)\widehat\Delta=g+\delta g$ with $\norm{\delta H}=O(u)\norm H$,
$\norm{\delta g}=O(u)\norm g$ and constants independent of $P$, where $u$ is the unit
roundoff. This is the provable numerical stability requested.

\section{Summary of what is established}\label{sec:summary}

\begin{itemize}
\item \textbf{(i) Feasibility decision:} Phase~I (\S\ref{sec:phaseI}) solves
$\min\tau$ \eqref{eq:phaseIopt} to the sign of $\tau^\star$;
the \emph{strict} PLMI is feasible in the open box, $\mathcal F\cap\{\norm{\bm x}<R\}
\neq\emptyset$, iff $\tau^\star<0$. When infeasible ($\tau^\star\ge0$) it returns an
\emph{explicit dual certificate} $\{Z_k\succeq0\}$, $\sum_k\tr Z_k=1$, satisfying the
stationarity \eqref{eq:dual} with $\beta\le0$; this gives $\sum_k\langle Z_k,\mathcal
A_k(\bm x)\rangle=\beta\le0$ for all $\bm x$, which (with $Z_k\succeq0$) is
incompatible with a strict solution $\mathcal A_k(\bm x)\succ0$ unless all
$\langle Z_k,\mathcal A_k\rangle=0$, the marginal boundary case
(Theorem~\ref{thm:conv}).
\item \textbf{(ii) Periodic solution:} when feasible, the centred Phase~II iterate is
a strictly feasible $P$-periodic $X(0),\dots,X(P-1)$ (Theorem~\ref{thm:conv}).
\item \textbf{(iii) No $nP\times nP$ lift; structure preserving:} the only large
object touched is the SPD block-cyclic-tridiagonal Hessian, solved in $O(Pm^3)$ by
\eqref{eq:thomas}--\eqref{eq:woodbury2} (Lemma~\ref{lem:exact}); residual algebra is
done in periodic Schur coordinates (\S\ref{sec:schur}).
\item \textbf{Convergence:} $O(\sqrt{nP}\log\frac1\varepsilon)$ Newton steps
(Theorem~\ref{thm:conv}).
\item \textbf{Per-iteration complexity:} $O\!\big(P(m^2n^2+mn^3+n^3)\big)$ flops,
$O(P(m+n^2))$ storage, both linear in $P$; total
$O\!\big(P^{3/2}n^{1/2}(m^2n^2+mn^3+n^3)\log\frac1\varepsilon\big)$
(Theorem~\ref{thm:cost}), strictly better in $P$ than the lifted $O(m^3P^3)$ solver.
\item \textbf{Stability:} orthogonal periodic Schur transformations plus pivoted SPD
block factorisation give backward stability with $P$-independent constants
(\S\ref{sec:schur}).
\end{itemize}

\end{document}
